% Human source transcription, not native author TeX.
% https://ocw.mit.edu/courses/6-441-information-theory-spring-2016/70ca499b4657031ee80dedff6f4a5a22_MIT6_441S16_chapter_12.pdf
% 第12章§12.1，Theorem 12.1及其两个方向、严格与非严格阈值的证明，印刷页131--132。明确收录有限支持实随机变量这一实例；不收入仅给证明纲要的Sanov定理。

\paragraph{Notation (editorial, following the source).}
$D(Q\|P)$ denotes relative entropy, with natural logarithms. The binary relative entropy is $d(q\|p)=q\log(q/p)+(1-q)\log((1-q)/(1-p))$, and $h(q)=-q\log q-(1-q)\log(1-q)$. Write $X^n=(X_1,\ldots,X_n)$ and $P_{X^n}=P^n$. The argument below is recorded for a real random variable with finite support; all expectations used in the minimizations are then well defined. This is an explicitly restricted instance of the source's Theorem 12.1.

\begin{theorem}[Large-deviation exponents, Theorem 12.1]
Let $X^n\stackrel{\mathrm{i.i.d.}}{\sim}P$. Then for any $\gamma\in\R$,
\begin{align}
\lim_{n\to\infty}\frac1n\log\frac1{P[\frac1n\sum_{k=1}^n X_k>\gamma]}
&=\inf_{Q:\E_Q[X]>\gamma}D(Q\|P),\label{ld:strict}\\
\lim_{n\to\infty}\frac1n\log\frac1{P[\frac1n\sum_{k=1}^n X_k\ge\gamma]}
&=\inf_{Q:\E_Q[X]\ge\gamma}D(Q\|P).\label{ld:closed}
\end{align}
\end{theorem}
\begin{proof}
We first prove \eqref{ld:strict}. Set $P[E_n]=P[\frac1n\sum_{k=1}^nX_k>\gamma]$.

\textbf{Lower Bound on $P[E_n]$:} Fix a $Q$ such that $\E_Q[X]>\gamma$. Let $X^n$ be iid. Then by WLLN,
\[
Q[E_n]=Q\left[\sum_{k=1}^nX_k>n\gamma\right]=1-o(1).
\]
Now the data processing inequality gives
\[
d(Q[E_n]\|P[E_n])\le D(Q_{X^n}\|P_{X^n})=nD(Q\|P).
\]
And a lower bound for the binary divergence is
\[
d(Q[E_n]\|P[E_n])\ge-h(Q[E_n])+Q[E_n]\log\frac1{P[E_n]}.
\]
Combining the two bounds on $d(Q[E_n]\|P[E_n])$ gives
\begin{equation}\label{ld:change}
P[E_n]\ge\exp\left(\frac{-nD(Q\|P)-\log2}{Q[E_n]}\right).
\end{equation}
Optimizing over $Q$ to give the best bound:
\[
\limsup_{n\to\infty}\frac1n\log\frac1{P[E_n]}
\le\inf_{Q:\E_Q[X]>\gamma}D(Q\|P).
\]

\textbf{Upper Bound on $P[E_n]$:} The key observation is that given any $X$ and any event $E$, $P_X(E)>0$ can be expressed via the divergence between the conditional and unconditional distribution as:
\[
\log\frac1{P_X(E)}=D(P_{X\mid X\in E}\|P_X).
\]
Define $\widetilde P_{X^n}=P_{X^n\mid\sum X_i>n\gamma}$, under which $\sum X_i>n\gamma$ holds a.s. Then
\begin{equation}\label{ld:condition}
\log\frac1{P[E_n]}=D(\widetilde P_{X^n}\|P_{X^n})
\ge\inf_{Q_{X^n}:\E_Q[\sum X_i]>n\gamma}D(Q_{X^n}\|P_{X^n}).
\end{equation}
We know show that the last problem ``single-letterizes'', i.e. need to be solved only for $n=1$. Consider the following two steps:
\begin{align}
D(Q_{X^n}\|P_{X^n})&\ge\sum_{j=1}^nD(Q_{X_j}\|P)\label{ld:product}\\
&\ge nD(\overline Q\|P),\qquad
\overline Q\mathrel{:=}\frac1n\sum_{j=1}^nQ_{X_j},\label{ld:convex}
\end{align}
where the first step follows from Corollary 2.1 after noticing that $P_{X^n}=P^n$, and the second step is by convexity of divergence Theorem 4.1. From this argument we conclude that
\begin{align}
\inf_{Q_{X^n}:\E_Q[\sum X_i]>n\gamma}D(Q_{X^n}\|P_{X^n})
&=n\inf_{Q:\E_Q[X]>\gamma}D(Q\|P),\label{ld:single-strict}\\
\inf_{Q_{X^n}:\E_Q[\sum X_i]\ge n\gamma}D(Q_{X^n}\|P_{X^n})
&=n\inf_{Q:\E_Q[X]\ge\gamma}D(Q\|P).\label{ld:single-closed}
\end{align}
In particular, \eqref{ld:condition} and \eqref{ld:single-strict} imply the required lower bound in \eqref{ld:strict}.

Next we prove \eqref{ld:closed}. First, notice that the lower bound argument \eqref{ld:condition} applies equally well, so that for each $n$ we have
\[
\frac1n\log\frac1{P[\frac1n\sum_{k=1}^nX_k\ge\gamma]}
\ge\inf_{Q:\E_Q[X]\ge\gamma}D(Q\|P).
\]
To get a matching upper bound we consider two cases:
\begin{itemize}
\item \textbf{Case I:} $P[X>\gamma]=0$. If $P[X\ge\gamma]=0$, then both sides of \eqref{ld:closed} are $+\infty$. If $P[X=\gamma]>0$, then
\[
P\left[\sum X_k\ge n\gamma\right]
=P[X_1=\cdots=X_n=\gamma]=P[X=\gamma]^n.
\]
For the right-hand side, since $D(Q\|P)<\infty\Rightarrow Q\ll P\Rightarrow Q(X\le\gamma)=1$, the only possibility for $\E_Q[X]\ge\gamma$ is that $Q(X=\gamma)=1$, i.e., $Q=\delta_\gamma$. Then
\[
\inf_{\E_Q[X]\ge\gamma}D(Q\|P)=\log\frac1{P(X=\gamma)}.
\]
\item \textbf{Case II:} $P[X>\gamma]>0$. Since $P[\sum X_k\ge n\gamma]\ge P[\sum X_k>n\gamma]$, from \eqref{ld:strict} we know that
\[
\limsup_{n\to\infty}\frac1n\log\frac1{P[\frac1n\sum_{k=1}^nX_k\ge\gamma]}
\le\inf_{Q:\E_Q[X]>\gamma}D(Q\|P).
\]
We next show that in this case
\begin{equation}\label{ld:boundary}
\inf_{Q:\E_Q[X]>\gamma}D(Q\|P)=\inf_{Q:\E_Q[X]\ge\gamma}D(Q\|P).
\end{equation}
Indeed, let $\widetilde P=P_{X\mid X>\gamma}$ which is well defined since $P[X>\gamma]>0$. For any $Q$ such that $\E_Q[X]\ge\gamma$, set $\widetilde Q=\bar\eps Q+\eps\widetilde P$ satisfies $\E_{\widetilde Q}[X]>\gamma$. Then by convexity,
\[
D(\widetilde Q\|P)\le\bar\eps D(Q\|P)+\eps D(\widetilde P\|P)
=\bar\eps D(Q\|P)+\eps\log\frac1{P[X>\gamma]}.
\]
Sending $\eps\to0$, we conclude the proof of \eqref{ld:boundary}.
\end{itemize}
\end{proof}
